%=====================================================================
%  SULI Research Report Paper
%  Formatted per the AIP Style Manual guidelines excerpted in
%  FinalProject/Instructions.txt:
%    - 12 pt font, all pages numbered from the title/abstract page
%    - title in sentence case, placed ~1/3 down page 1
%    - single-paragraph abstract, double spaced, wider margins
%    - AIP heading scheme: I. PRINCIPAL / A. Sub / 1. / a.
%    - numbered TABLE (roman) and FIG. (arabic), each cited in text
%    - ACKNOWLEDGMENTS, then APPENDIX, then REFERENCES
%    - superscript numeral citations
%
%  COMPILE:  latexmk -pdf report.tex   (or Ctrl+Alt+B in VS Code)
%  Figures live in ../report_figures/ and ../ (see \graphicspath).
%=====================================================================
\documentclass[12pt]{article}

\usepackage[T1]{fontenc}
\usepackage[utf8]{inputenc}
\usepackage{lmodern}
\usepackage[margin=1in]{geometry}
\usepackage{graphicx}
\usepackage{booktabs}
\usepackage{amsmath}
\usepackage{amssymb}
\usepackage{xcolor}
\usepackage{caption}
\usepackage{float}
\usepackage{enumitem}
\usepackage{microtype}
\usepackage{setspace}
\usepackage{changepage}
\usepackage{titlesec}
\usepackage{indentfirst}
\usepackage[hidelinks]{hyperref}

\graphicspath{{../report_figures/}{../FinalProject/Figures/roi_shape_comparison/}{../}}

% --- AIP heading numbering: I. / A. / 1. / a. ---
\renewcommand{\thesection}{\Roman{section}}
\renewcommand{\thesubsection}{\Alph{subsection}}
\renewcommand{\thesubsubsection}{\arabic{subsubsection}}
\titleformat{\section}[block]{\centering\normalfont\bfseries}{\thesection.}{0.6em}{\MakeUppercase}
\titleformat{\subsection}[block]{\normalfont\bfseries}{\thesubsection.}{0.5em}{}
\titleformat{\subsubsection}[block]{\normalfont\itshape}{\thesubsubsection.}{0.5em}{}
\titlespacing*{\section}{0pt}{1.4em}{0.8em}

% --- TABLE (roman) / FIG. (arabic), captions labelled per AIP ---
\renewcommand{\thetable}{\Roman{table}}
\captionsetup[table]{name=TABLE,labelsep=period,font=small,labelfont=bf,justification=raggedright,singlelinecheck=false}
\captionsetup[figure]{name=FIG.,labelsep=period,font=small,labelfont=bf}

% --- page numbers on every page (default plain footer) ---
\pagestyle{plain}

% --- helpers ---
% Inline red flag for a gap to resolve before final submission.
\newcommand{\todo}[1]{\textcolor{red}{\textbf{[TODO: #1]}}}
% Superscript citation.
\newcommand{\cq}[1]{\textsuperscript{#1}}
% Missing-value marker for tables (AIP: raised dots, not dashes).
\newcommand{\na}{\ensuremath{\ldots}}
\newcommand{\um}{\ensuremath{\mu}m}

\begin{document}

%=====================================================================
% TITLE + ABSTRACT PAGE  (title ~1/3 down, numbered page 1)
%=====================================================================
\thispagestyle{plain}
\vspace*{1.8in}
\begin{center}
{\large\bfseries High-resolution structural mapping of perovskite thin films:
resolving (111) and (001) microstructure via automated nano-XRD analysis and a
CLI-based computer vision pipeline\par}
\vspace{1.5em}
{\normalsize Takaji Robson\par}
\vspace{0.4em}
{\small X-ray Science Division, Argonne National Laboratory\\
In Situ Nanoprobe (ISN), Advanced Photon Source, beamline 19-ID\par}
\vspace{0.4em}
{\small Science Undergraduate Laboratory Internship (SULI), 2026\par}
\end{center}

\vspace{1.5em}
% --- Abstract: one paragraph, single spaced, wider margins, no math/figs/tables ---
\begin{adjustwidth}{0.5in}{0.5in}
\noindent\textbf{ABSTRACT.}\ Nano-focused X-ray diffraction (nano-XRD) at the
Advanced Photon Source (APS) In Situ Nanoprobe (beamline 19-ID) resolves crystal
grain size and orientation in semiconductor thin films, but each scan generates
on the order of $10^5$ diffraction frames, creating a bottleneck for rapid
analysis. We developed xrd-app, a command-line computer-vision pipeline that
processes raw HDF5 nano-XRD data through frame-by-frame Bragg peak detection and
spatial linking of neighboring detections into grain-level structural features.
An automated algorithm-discovery system (CVEvolve) tunes the peak detector, and
spatial binning raises the per-position signal-to-noise ratio approximately
linearly with bin size, improving detection of weaker peaks. A contouring stage
based on Union-Find clustering and Gaussian-profile filtering merges scattered
detections into coherent grain outlines. Applied to six focused scans across
three perovskite-halide solar cells, the workflow shows that samples processed
outside a glove box adopt a stronger preferred (001) orientation relative to
(111), while a deionized-water (DI) treatment increases the average size of
(001)-oriented grains.
\end{adjustwidth}

\vspace{1em}

%=====================================================================
\section{Introduction}
%=====================================================================
Metal-halide perovskite thin films are leading candidates for next-generation
photovoltaics because they combine high power-conversion efficiency with
low-cost, solution-based fabrication. Device performance is strongly coupled to
microstructure: the size, orientation, and spatial distribution of crystalline
grains govern charge transport and defect density. Synchrotron nano-focused
X-ray diffraction (nano-XRD) is uniquely suited to probing this microstructure,
because a nano-focused beam can map the local crystallographic orientation and
grain size across a working device with sub-micrometer spatial
resolution.\cq{1}

At the APS In Situ Nanoprobe (beamline 19-ID), a single nano-XRD raster scan
records a two-dimensional charge-coupled-device (CCD) diffraction frame at every
spatial position, so a full scan comprises on the order of $10^5$ frames stored
as compressed HDF5. Extracting grain-level information from such a dataset by
hand is impractical, and the volume of data, tens to hundreds of gigabytes per
scan, creates a computational bottleneck for rapid, reproducible analysis.

The purpose of this work is to develop and validate an automated computer-vision
pipeline, xrd-app, that converts raw nano-XRD frames into physically meaningful
grain maps with minimal human intervention, and to apply it to a set of
perovskite-halide devices to test whether processing conditions leave a
measurable fingerprint on grain microstructure. This paper is organized as
follows. Section~II describes the data and the three processing stages: binning,
peak detection with automated detector optimization, and grain contouring.
Section~III reports the effect of each stage on data quality and then
presents the grain-level comparison across devices. Section~IV summarizes the
findings and outlines future work.

%=====================================================================
\section{Methods}
%=====================================================================

\subsection{Nano-XRD data acquisition and structure}
Measurements were performed at APS beamline 19-ID using a 15\,keV nano-focused
beam and a CCD area detector with a 75\,\um\ pixel pitch. Each raster position
yields one detector frame; the pixel-to-$2\theta$ mapping is fixed by a
calibration image, and the spatial position of each frame is recorded by
interferometer readout. Reflections of interest for the perovskite films and the
indium tin oxide (ITO) substrate, including PbI$_2$, (001), (011), (111),
(002), and (012)/(112), appear in known radial ($2\theta$) bands, which define
where genuine Bragg peaks are expected. A representative scan (Scan\_0203)
contains 25{,}170 occupied spatial positions; the six focused scans analyzed
here range from roughly 80 to 200\,GB each. Basic instrument and dataset
parameters are collected in Table~\ref{tab:dataset}.

\begin{table}[H]
\centering
\caption{Instrument and dataset parameters for the reference scan (Scan\_0203).}
\label{tab:dataset}
\small
\begin{tabular}{@{}ll@{}}
\toprule
Parameter & Value\\
\midrule
Beamline & APS 19-ID (In Situ Nanoprobe)\\
Photon energy & 15\,keV\\
Detector pixel pitch & 75\,\um\\
Frame size & $1062\times1028$ px (full); $256\times256$ px (region of interest)\\
Reflections & PbI$_2$, (001), (011), (111), (002), ITO, (012), (112)\\
Unbinned grid & 156 rows $\times$ 221 cols $=$ 34{,}476 cells; 25{,}170 occupied\\
$5\times5$ binned grid & 32 rows $\times$ 45 cols $=$ 1{,}440 cells; 1{,}188 occupied\\
\bottomrule
\end{tabular}
\end{table}

\subsection{Spatial binning}
Because a single-position frame often contains only a weak diffraction signal,
the pipeline can sum $N\times N$ neighboring frames prior to detection. Binning
reduces the spatial grid size (for example, from $151\times145$ to $51\times49$)
but raises the per-position signal-to-noise ratio (SNR) approximately linearly
with the bin dimension $N$, because coherent signal accumulates faster than
random noise. Binning also stabilizes grain outlines by averaging over the
small, systematic position offsets introduced by the serpentine (S-shaped)
raster trajectory; in the unbinned case these offsets can appear as holes or
horizontal breaks in the reconstructed grain maps. To bin on a physically
faithful lattice, the pipeline reads the true interferometer coordinates, snaps
each raw frame to its nearest node on a square-pixel grid whose origin is the
minimum $(x,y)$ of the scan, and uses the median inter-point distance to set the
node spacing. Figure~\ref{fig:binning} illustrates the concept.

\begin{figure}[H]
\centering
\includegraphics[width=\linewidth]{binning_faithful_grid}
\caption{Faithful spatial binning. The true stage positions (left) follow the
S-shaped raster; each frame is snapped to the nearest node of a square grid
(center), where robust min/max limits clip the positions that jut past the edges
and a ``2'' marks nodes where two frames coincide; neighboring nodes are then
summed into $3\times3$ bins (right), trading spatial resolution for higher
per-position signal-to-noise.}
\label{fig:binning}
\end{figure}

\subsection{Peak detection and automated detector optimization}
Bragg peaks are detected frame by frame within the expected $2\theta$ bands.
Rather than hand-tuning a single detector, we used CVEvolve, an automated
algorithm-discovery system that iteratively generates and tunes candidate
detection algorithms and scores each against ground-truth
annotations.\cq{2} Candidate detectors were scored with the mean $F_2$
statistic (a recall-weighted harmonic mean, $\beta\approx2$) rather than $F_1$,
because the ground truth is derived from coarser $3\times3$ annotations mapped to
the central frame and therefore undercounts the peaks visible at full unbinned
sensitivity. Weighting recall roughly four times precision avoids penalizing
genuine detections that the sparse ground truth happens to omit. Matches were
accepted within a 40\,px tolerance.

\subsection{Grain contouring}
A detection in a single spatial bin (a \emph{peak}) may be noise; a physically
real feature is a peak that persists across neighboring positions. The
contouring stage links per-position detections into \emph{grains} by combining a
Union-Find spatial clustering step with a Gaussian-profile filter that retains
clusters having a bright, monotonic center and rejects fragmented or noisy
detections. Each retained grain is characterized by its footprint (number of
bins), azimuthal orientation $\chi$, intensity, and radial breadth
($\Delta 2\theta$). This step permits a less restrictive per-frame detection
threshold, because it assumes grain persistence across neighboring pixels, at
the cost of imposing an effective minimum grain size. Figures~\ref{fig:contour}
and~\ref{fig:contourslice} show a representative contoured grain and its
intensity cross-section.

\begin{figure}[H]
\centering
\begin{minipage}{0.48\linewidth}\centering
  \includegraphics[width=\linewidth]{contour_example}
  \captionof{figure}{A contoured grain reconstructed by Union-Find linking of
  neighboring per-position detections.}
  \label{fig:contour}
\end{minipage}\hfill
\begin{minipage}{0.48\linewidth}\centering
  \includegraphics[width=\linewidth]{contour_slice}
  \captionof{figure}{Intensity cross-section through a grain; the
  Gaussian-profile filter keeps bright, monotonic centers.}
  \label{fig:contourslice}
\end{minipage}
\end{figure}

%=====================================================================
\section{Results and discussion}
%=====================================================================

\subsection{Signal-to-noise and feature yield versus bin size}
Table~\ref{tab:counts} lists the number of occupied bins, raw peaks, and
retained grains as a function of bin size for Scan\_0203. Both the total number
of raw peaks and the number of retained grains \emph{decrease} with bin size,
because the bin count falls as $\sim1/N^2$, faster than the per-bin yield rises.
However, the quality of each detection improves sharply: as shown in
Table~\ref{tab:snr}, the median peak SNR rises from about 5.3 at $1\times1$
(barely above threshold) to roughly 17--24 when binned, and the median grain
intensity rises tenfold. Figure~\ref{fig:snrhist} shows the full SNR
distribution shifting to higher values with bin size. The apparent paradox of
fewer but stronger
detections arises because most bins contain only background and a
sub-bin-sized grain does not fill a larger bin, so summed signal grows more
slowly than $N^2$ while noise grows as $N$.

\begin{table}[H]
\centering
\caption{Occupied bins, raw peaks, and retained grains versus bin size
(Scan\_0203; detector \texttt{tophat\_band\_adaptive\_snr}, SNR threshold 4,
linking tolerance 5\,px).}
\label{tab:counts}
\small
\begin{tabular}{@{}lrrrl@{}}
\toprule
Bin size & Occupied bins & Raw peaks & Retained grains & Spatial resolution\\
\midrule
$1\times1$ & 25{,}170 & 186{,}063 & 1{,}825 & $156\times221$\\
$2\times2$ & 7{,}231  & 16{,}915  & 310   & $78\times111$\\
$3\times3$ & 3{,}253  & 9{,}124   & 213   & $52\times74$\\
$4\times4$ & 1{,}859  & 6{,}125   & 199   & $39\times56$\\
$5\times5$ & 1{,}188  & 3{,}889   & 144   & $32\times45$\\
\bottomrule
\end{tabular}
\end{table}

\begin{table}[H]
\centering
\caption{Signal-to-noise ratio and intensity at real detections rise with
binning even as total counts fall (Scan\_0203).}
\label{tab:snr}
\small
\begin{tabular}{@{}lll@{}}
\toprule
Quantity & $1\times1$ & Binned\\
\midrule
Median peak SNR & 5.3 & 17--24\\
90th-percentile peak SNR & 27 & 357\\
Median grain intensity (arb.) & 6 & 64\\
\bottomrule
\end{tabular}
\end{table}

\begin{figure}[H]
\centering
\includegraphics[width=0.72\linewidth]{snr_hist_by_binsize}
\caption{Distribution of detection SNR at each bin size; the distribution shifts
to higher SNR as bin size grows.}
\label{fig:snrhist}
\end{figure}

The file size of the pre-built binned HDF5 does not fall as fast as the bin count
(10\,GB at $1\times1$ versus 971\,MB at $5\times5$), because each stored bin
retains a full-resolution detector image regardless of bin size. Building a
binned file is the throughput bottleneck, running roughly one to two orders of
magnitude slower than detection itself. Taken together, a $3\times3$ bin is a
practical sweet spot: the SNR gain has largely saturated and the horizontal
fragmentation is gone, while the spatial resolution ($52\times74$) remains high
enough to resolve device-scale structure.

\subsection{Outline robustness relative to fixed detector-ROI thresholding}
A conventional alternative to linked peak detection is to place a fixed detector
region of interest (ROI) around a known Bragg spot, integrate that aperture at
every scan position, and define the grain as the connected region above a chosen
SNR limit. We compared this approach with xrd-app for Scan\_0203 using a
$7\times7$-pixel signal aperture, a surrounding $17\times17$-pixel background
window, and the same $3\times3$ spatial bins. As a common reference we built a
\emph{territorial} map: the same Union-Find linking of per-position detections,
but grouped on each frame's real, imperfect interferometer position rather than
on the regularized square grid. Because it uses the true stage coordinates and
no outline threshold of its own, it provides a method-independent footprint
against which both the xrd-app and fixed-ROI outlines can be scored. Features
were matched across the regular and territorial maps by reflection, detector
position, and overlap among their underlying raw frames.
Of the credible matches, three representative cases were selected objectively at
the 90th, 50th, and 10th percentiles of elliptical-Gaussian fit quality
($R^2$), rather than by visual preference.

Figure~\ref{fig:roicomparison} shows the territorial footprint, xrd-app outline,
fixed-ROI outline at SNR 4, and an edge profile for each case. In the final
column the two SNR traces are computed by different conventions and are not
directly comparable in absolute height: the orange xrd-app trace is the detector
SNR after radial-background subtraction and top-hat filtering, evaluated at the
peak wherever it lands on the detector in each frame, whereas the open-blue trace
is plain aperture photometry in one static detector box. The panel is therefore
meant to be read for the \emph{shape} of the falloff and the fact that the
peak-following trace stays high across the grain while the fixed aperture decays,
not as a claim that one records more signal. The Gaussian curve is a fit used
only to display profile breadth; it is not an acceptance step in the production
pipeline, which instead requires a non-flat, generally decreasing intensity
profile after adaptive peak detection and spatial linking. This distinction
matters for the poor case: its spatial
profile is visibly non-Gaussian ($R^2=0.30$), but local detector continuity still
recovers a coherent footprint with an intersection-over-union (IoU) of 0.61
against the territorial reference. The fixed ROI reaches at most 0.52 for this
feature over the tested thresholds and only 0.32 at SNR 4. For the ideal and
median cases, xrd-app gives IoU values of 0.43 and 0.49, respectively. A tuned
fixed ROI approaches these values, showing that simple thresholding remains
adequate for compact, high-quality profiles when the correct cutoff is already
known.

\begin{figure}[p]
\centering
\includegraphics[width=\textwidth]{roi_vs_xrdapp_examples}
\caption{Fixed detector-ROI thresholding compared with xrd-app for matched
Scan\_0203 features. Rows show ideal, median, and poor Gaussian-fit cases (90th,
50th, and 10th percentiles of fit $R^2$). Columns: the territorial reference
(Union-Find on the real, imperfect stage positions), the $3\times3$ xrd-app
linked shape, the fixed-ROI region at SNR 4, and the SNR dropoff along the fitted
major axis. Orange points are xrd-app detections that follow local Bragg-peak
movement (background-subtracted, top-hat SNR); open blue points use one fixed
aperture. The two use different SNR conventions, so compare the falloff shape,
not the absolute height. The poor case shows xrd-app keeps a reproducible outline
even when the profile is not a single Gaussian.}
\label{fig:roicomparison}
\end{figure}

The principal weakness of fixed-ROI segmentation is therefore not that every
chosen cutoff fails, but that the required cutoff and resulting area are
feature-dependent. Figure~\ref{fig:roisensitivity} sweeps the conventional SNR
limit from 2 to 8. The inferred area changes from 13 to 35 bins for the ideal
case, from 0 to 48 bins for the median case, and from 65 to 269 bins for the poor
case. The median feature disappears completely at SNR 7 even though xrd-app
retains a linked 16-bin shape. In contrast, the saved xrd-app outline needs no
feature-specific threshold after detection: a common detector SNR threshold, local
position matching, and a profile-consistency test make the boundary less sensitive
to noisy or missing positions while still allowing imperfect physical shapes.

\begin{figure}[H]
\centering
\includegraphics[width=\textwidth]{roi_threshold_sensitivity}
\caption{Sensitivity of conventional fixed-ROI outlines to the selected SNR
cutoff. Blue curves give frame-footprint IoU with the independently binned
territorial reference; stars mark the best cutoff for each feature. Orange lines
give the IoU of the saved xrd-app shape, which is unchanged throughout the
conventional threshold sweep. The conventional area varies by factors of
approximately 2.7--4.1 and can vanish entirely for a moderate-SNR feature.}
\label{fig:roisensitivity}
\end{figure}

\subsection{Device-level grain statistics}
Applying the full workflow to six focused scans across three perovskite-halide
devices, distinguished by whether the film was processed in a glove box (GB) and
by a deionized-water (DI) treatment, reveals reproducible microstructural
differences. Figure~\ref{fig:gb} compares the glove-box and no-glove-box
conditions. Count dominance and size dominance behave differently: (001) remains
the more numerous orientation in every device, but the two orientations exchange
which one grows larger. In the no-glove-box devices the (001):(111) grain-size
ratio falls markedly (from $\sim$3.4 to $\sim$0.9) and the largest (111) grain
roughly doubles relative to the overall device area. In other words, processing
without glove-box control preserves the prevalence of (001) grains by number
while substantially enhancing the growth of individual (111) grains, so the
larger grains become predominantly (111).

Figure~\ref{fig:treat} compares the treatment conditions. The DI treatment
increases both the average size and the total area share of (001)-oriented
grains: the summed (001):(111) area ratio rises from about 2.9 to 4.6, while the
mean (001) grain size increases and the (111) population is relatively unchanged.
As a concrete example of the per-reflection output, for Scan\_0182 the pipeline
detected 78 (001) grains and 35 (111) grains on the $3\times3$ binned grid, and
208 (001) versus 104 (111) grains when reconstructed on the true-position grid.

\begin{figure}[H]
\centering
\begin{minipage}{0.48\linewidth}\centering
  \includegraphics[width=\linewidth]{conclusion1_GB}
  \captionof{figure}{Glove-box vs.\ no-glove-box: (001):(111) grain statistics.}
  \label{fig:gb}
\end{minipage}\hfill
\begin{minipage}{0.48\linewidth}\centering
  \includegraphics[width=\linewidth]{conclusion2_5pctDI}
  \captionof{figure}{Effect of the deionized-water (DI) treatment on (001) grain
  size and area share.}
  \label{fig:treat}
\end{minipage}
\end{figure}

%=====================================================================
\section{Conclusions and future work}
%=====================================================================
We developed xrd-app, an automated command-line computer-vision pipeline that
turns raw nano-XRD frames into grain-level structural maps, and validated it on
perovskite-halide solar cells measured at APS 19-ID. Spatial binning raises the
per-position signal-to-noise ratio approximately linearly with bin size, with a
practical optimum near $3\times3$, and a Union-Find plus Gaussian-profile
contouring stage consolidates scattered detections into physically meaningful
grains while tolerating a less restrictive detection threshold. Against a
true-position territorial reference, this linked approach stays effective for
nonideal profiles where a fixed ROI needs feature-specific tuning and can change
the inferred area severalfold. Across three devices, the workflow shows that
processing outside a glove box strongly enhances (111) grain growth relative to
(001), and that a deionized-water (DI) treatment increases the size and area share
of (001) grains, correlations obtained automatically from datasets impractical to
analyze by hand.

The workflow is most effective for samples with predictable, repeated Bragg
angles: once the expected reflections are known, feature finding is fully
automated across any number of scans and the grain maps are directly comparable
device to device. Even when the angles are unknown, binning, peak detection, and
Union-Find linking remain reproducible, scalable, and high-confidence in any case
where a minimum grain size and an intensity dropoff can be assumed. The automation
matters most for fine-grained films: when a device holds many small or
low-intensity grains, their reflections are lost in the noise gradient of the
fully summed image and cannot be placed by eye, so the pipeline scans every
position and recovers each peak (Fig.~\ref{fig:grain-scale}, left). For
coarse-grained devices a few dominant grains stand out in the summed image and can
be picked by hand (right), so automation is optional.

\begin{figure}[H]
\centering
\includegraphics[width=0.82\linewidth]{small_vs_large_grain}
\caption{Left: every (001) grain across Scan\_0179 (true-position territorial
reference); many faint features are invisible in a summed image and recovered
only by scanning each position. Right: a single large Scan\_0032 grain, obvious in
a fixed ROI on the summed image.}
\label{fig:grain-scale}
\end{figure}

Future work will (1) apply the same workflow to additional device types, (2)
correlate the nano-XRD grain maps with spatially aligned X-ray fluorescence
signatures to connect microstructure with elemental composition, and (3) extend
the analysis to rocking-curve series in order to reconstruct three-dimensional
spatial maps of the identified grains.

%=====================================================================
\section*{Acknowledgments}
\addcontentsline{toc}{section}{Acknowledgments}
%=====================================================================
This work was performed under the U.S. Department of Energy (DOE) Science
Undergraduate Laboratory Internship (SULI) program, sponsored by the DOE Office
of Science, Office of Workforce Development for Teachers and Scientists. The
author thanks his mentors Yanqi Grace Luo and Pedro Mercado Lozano for their
guidance throughout the project, and the staff of the X-ray Science Division and
the In Situ Nanoprobe at the Advanced Photon Source. This research used
resources of the Advanced Photon Source, a U.S. Department of Energy (DOE)
Office of Science user facility operated for the DOE Office of Science by Argonne
National Laboratory under Contract No. DE-AC02-06CH11357.

%=====================================================================
\section*{References}
\addcontentsline{toc}{section}{References}
%=====================================================================
\begin{enumerate}[label={\arabic*.},leftmargin=2em,itemsep=2pt]
\item M. E. Stuckelberger, T. Nietzold, B. M. West, Y. Luo, X. Li, J. Werner,
\emph{et al.}, ``Effects of X-rays on perovskite solar cells,'' J. Phys. Chem.
C \textbf{124}, 17949--17956 (2020).
\item M. Du, X. Yin, Y. Luo, D. Beniwal, S. Tang, H. Sharma, and M. J. Cherukara,
``CVEvolve: Autonomous algorithm discovery for unstructured scientific data
processing,'' arXiv:2605.11359 (2026).
\item Y. Li, Z. Zhang, H. Zhang, Y. Liu, H. Gao, and Y. Mao, ``Preparation of
(111) facet-dominated perovskite films for highly efficient and stable
perovskite solar cells,'' ACS Appl. Mater. Interfaces \textbf{17},
48231--48239 (2025).
\end{enumerate}

%=====================================================================
\section*{Appendix: Largest features, manual ROI versus automated detection}
\addcontentsline{toc}{section}{Appendix}
%=====================================================================

\begin{figure}[htbp]
\centering
\includegraphics[width=0.68\linewidth]{appendix_roi_vs_xrdapp}
\caption{The three largest features in Scan\_0032. Left: a manual detector ROI on
the fully summed diffraction image, giving total counts inside a fixed box at
every spatial bin; the feature sits in a bright, serpentine-striped background the
fixed box cannot separate from signal. Right: the same feature recovered by
xrd-app at $3\times3$ binning, plotted at its true grid location, isolated to a
clean footprint (white where no peak was linked). Both agree on location, but only
the automated map reports size and shape without manual thresholding.}
\label{fig:appendix-roi}
\end{figure}

\begin{figure}[htbp]
\centering
\includegraphics[width=\linewidth]{appendix_xrdapp_shapes}
\caption{Largest features in Scans 179 and 182, each shown as its linked
$3\times3$ footprint on a white background (the Shape/Verify view). No manual ROIs
were drawn here; the workflow finds and characterizes every feature without
operator input, so this view scales to any scan.}
\label{fig:appendix-shapes}
\end{figure}

\end{document}
